\documentclass{article} % This command is used to set the type of
% document you are working on such as an article, book, or presenation

\usepackage{geometry} % This package allows the editing of the page layout
\usepackage{amsmath}  % This package allows the use of a large range
% of mathematical formula, commands, and symbols
\usepackage{graphicx}  % This package allows the importing of images
\usepackage{amssymb}

\newcommand{\question}[2][]{
  \begin{flushleft}
    \textbf{Question #1}: \textit{#2}

\end{flushleft}}
\newcommand{\sol}{\textbf{Solution}:} %Use if you want a boldface solution line
\newcommand{\maketitletwo}[2][]{
  \begin{center}
    \Large{\textbf{Assignment #1}

    Topology} % Name of course here
    \vspace{5pt}

    \normalsize{Nimrod Ohayon Rozanes  % Your name here

    \today}        % Change to due date if preferred
    \vspace{15pt}

\end{center}}
\begin{document}
\maketitletwo[4]  % Optional argument is assignment number
%Keep a blank space between maketitletwo and \question[1]

\question[1]{}

By definition
$$
\text{Int}(Y) = \bigcup U_i
$$
For $\{U_i\}$ being the set of all open sets for which $U_i \subset
Y$. Then we want to show that

\[
  \bigcup U_i = \{x \vert \forall x \in X, \exists
  \text{nbhd}(x) \subset Y\}
\]

Since $\text{nbhd}(x) \subset Y$ and a neighborhood is open, then for
all $x, \exists \text{nbhd}(x) \in \{U_i\}$. Take the union of all
such neighborhoods, yeilding that

$$
\{x \vert \forall x \in X, \exists
\text{nbhd}(x) \subset Y\} \subset \bigcup \text{nbhd}(x) \subset \bigcup U_i
$$

Now we want to show that

\[
  \bigcup U_i \subset \{x \vert \forall x \in X, \exists
  \text{nbhd}(x) \subset Y\}
\]

Take any $y \in \bigcup U_i$, then $\exists i \in I, y \in U_i$, by
definition, $U_i$ is a neighborhood of $y$ which is a subset of $Y$.
Therefore $\forall y \in \bigcup U_i$, it is the case that

\[
  y  \in \{x \vert \forall x \in X, \exists \text{nbhd}(x) \subset Y\}
\]

as such
\[
  \bigcup U_i = \{x \vert \forall x \in X, \exists \text{nbhd}(x) \subset Y\}
\]

\question[2]{}
\begin{itemize}
  \item $\text{Int}(Y) = X \setminus \overline{X \setminus Y}$\\
    Consider that
    \[
      X \setminus \overline{X\setminus Y} = X \setminus \bigcap C_i =
      \bigcup X \setminus C_i
    \]
    for $\{C_i\}$ being the set of all sets such that $X\setminus Y
    \subset C_i$ and $X \setminus C_i$ is open. Further, consider that
    \[
      X \setminus Y \subset C_i \iff X \subset C_i
      \cup Y \iff X \setminus C_i \subset Y
    \]

    Clearly, $X \setminus \overline{X \setminus Y}$ is the union of
    all open sets which are contained inside of $Y$, which is the
    definition of $\text{Int}(Y)$
  \item $\overline Y = X \setminus \text{Int}(X\setminus Y)$\\
    Consider that

    \[
      X \setminus \text{Int}(X\setminus Y) = X \setminus \bigcup U_i
      = \bigcap X \setminus U_i
    \]

    for $\{U_i\}$ the set of all sets which are open and contained in $X
    \setminus Y$. Since $U_i$ is open, $X \setminus U_i$ is closed, further,
    \[
      U_i \subset X \setminus  Y \iff U_i \cup Y \subset X
      \iff Y \subset X \setminus U_i
    \]
    Therefore, $X \setminus \text{Int}(X \setminus Y)$ is the
    intersection of all closed sets which contain $Y$, which is the
    defintion of the closure of $Y$
\end{itemize}
\question[3]{}
\begin{itemize}
  \item $\overline{Y \cup Z} = \overline{Y} \cup \overline{Z}$\\
    First we show $\overline{Y}  \cup \overline{Z} \subset
    \overline{Y \cup Z}$: We have that $Y \subset Y \cup Z \implies
    \overline Y \subset \overline{Y \cup Z}$ and that
    $Z \subset Y \cup Z \implies \overline{Z} \subset \overline{Y
    \cup Z}$, then $\overline Y \cup \overline Z \subset \overline{Y
    \cup Z}$. To show that $\overline{Y \cup Z} \subset
    \overline Y \cup  \overline Z$: since $Y \subset \overline Y$ and
    $Z \subset \overline Z$, we have that $Y \cup Z \subset \overline
    Y \cup  \overline Z \implies \overline{Y \cup Z} \subset
    \overline{\overline{Y} \cup \overline{Z}}$. Since $\overline Y$
    and $\overline Z$ are
    closed, so is $\overline Y \cup \overline Z$ as such
    $\overline{\overline Y \cup \overline Z} = \overline Y \cup
    \overline Z$, therefore, $\overline{Y \cup Z} \subset \overline Y
    \cup \overline Z$ as desired.

  \item $\overline{Y \cap Z}  \subset \overline Y \cap \overline Z$

    \begin{align*}
      Y \cap Z \subset Y \text{ and } Y \cap Z \subset Z &\implies
      \overline{(Y \cap Z)} \subset \overline Y \text{ and } \\
      &\implies \overline{Y \cap Z} \subset \overline Y \cap \overline Z
    \end{align*}
    The reverse inclusion does not hold in general: take $X = \mathbf{R}$
    with the standard topology, $Y = (0,1)$ and $Z = (1,2)$. Then $Y \cap
    Z = \emptyset$, so $\overline{Y \cap Z} = \emptyset$, but
    \[
      \overline Y \cap \overline Z = [0,1] \cap [1,2] = \{1\}
      \not\subset \emptyset = \overline{Y \cap Z}
    \]
  \item $\text{Int}(Y) \cup \text{Int}(Z) \subset \text{Int}(Y \cup Z)$\\
    \[
      Y \subset Y \cup Z \text{ and } Z \subset Y \cup Z \text{ and }
      \text{Int}(Y) \subset Y \text{ and } \text{Int}(Z) \subset Z
    \]
    then
    \[
      \text{Int}(Y) \cup \text{Int}(Z) \subset Y \cup Z
    \]
    then
    \[
      \text{Int}(\text{Int}(Y) \cup  \text{Int}(Z)) \subset \text{Int}(Y \cup Z)
    \]
    Since $\text{Int}(Y)$ and $\text{Int}(Z)$ are by definition open,
    their union is as well, therefore, $\text{Int}(\text{Int}(Y) \cup
    \text{Int}(Z)) = \text{Int}(Y) \cup \text{Int}(Z)$, as such:
    \[
      \text{Int}(Y) \cup  \text{Int}(Z) \subset \text{Int}(Y \cup Z)
    \]
    The reverse inclusion does not hold in general: take $X = \mathbf{R}$
    with the standard topology, $Y = [0,1]$ and $Z = [1,2]$. Then
    $\text{Int}(Y) \cup \text{Int}(Z) = (0,1) \cup (1,2)$, but
    \[
      \text{Int}(Y \cup Z) = \text{Int}([0,2]) = (0,2)
    \]
    so $1 \in \text{Int}(Y \cup Z)$ but $1 \notin \text{Int}(Y) \cup
    \text{Int}(Z)$.
  \item $\text{Int}(Y) \cap \text{Int}(Z) = \text{Int}(Y\cap Z)$ \\
    First we show that $\text{Int}(Y\cap Z) \subset \text{Int}(Y)
    \cap \text{Int}(Z)$:
    \begin{align*}
      Y \cap Z  \subset Y  \text{ and } Y \cap Z  \subset Z
      &\implies  \text{Int}(Y \cap Z) \subset \text{Int}(Y) \text{
      and } \text{Int}(Z) \\
      &\implies \text{Int}(Y \cap Z)\subset \text{Int}(Y) \cap \text{Int}(Z)
    \end{align*}
    Then also that $\text{Int}(Y) \cap \text{Int}(Z) \subset
    \text{Int}(Y \cap Z)$:
    \begin{align*}
      \text{Int}(Y) \subset Y \text{ and } \text{Int}(Z) \subset Z
      &\implies \text{Int}(Y) \cap \text{Int}(Z) \subset Y \cap Z \\
      &\implies \text{Int}(\text{Int}(Y) \cap \text{Int}(Z)) \subset
      \text{Int}A(Y \cap Z) \\
      &\implies \text{Int}(Y) \cap \text{Int}(Z) \subset \text{Int}(Y \cap Z)
    \end{align*}

\end{itemize}
\question[4]{}
\begin{itemize}
  \item $\partial \mathcal{I} = \{x \in \mathbb{R} \mid U \cap
      \mathcal{I} \neq \emptyset \text{ and } U \cap \mathcal{I}^c \neq
    \emptyset \text{ for every nbhd } U \text{ of } x\}$\\
    Take any $x \in \mathbb{R}$ and any neighborhood of it such that
    $x \in U = (a, b]$ for some $a, b\in \mathcal{R}$. If $x$ is
    irrational, that is $x \in \mathcal{I}$, then $x \in U \cap
    \mathcal{I} \neq \emptyset$ is trivial. Therefore, it suffcies to
    show that $U \cap \mathcal{I}^c \neq \emptyset$, that is, that
    for any neighborhood of a real number $x$ in the upper limit
    topology, there is a rational number included. Since $a,b \in
    \mathbb{R}$, and $a < b$ then by the archimedean property, there
    exists $r \in \mathbb{Q}$ such that $a < r < b$, and therefore $U
    \cap \mathcal{I}^c \neq \emptyset$ \\

    Similarily, if $x$ is rational, then $U \cap \mathcal{I}^c \neq
    \emptyset$ trivially, so it suffices to show that $U \cap
    \mathcal{I} \neq \emptyset$, to show that this is the case, take
    $x - a = \varepsilon$. Since $a < x$, $\varepsilon > 0$, so by the
    archimedean property there exists $n \in \mathbb{N}$ such that
    $\frac{\sqrt{2}}{n} < \varepsilon$. Let $y = x -
    \frac{\sqrt{2}}{n}$, then $a < y < x \le b$, so $y \in U$. Further,
    $y$ is irrational, since otherwise $\sqrt{2} = n(x - y)$ would be
    rational. Therefore $y \in U \cap \mathcal{I}$, and $U \cap
    \mathcal{I} \neq \emptyset$. \\

    Since every neighborhood of $x$ contains a basis element $(a, b]$
    containing $x$, every neighborhood of every $x \in \mathbb{R}$
    intersects both $\mathcal{I}$ and $\mathcal{I}^c$, as such
    \[
      \partial \mathcal{I} = \mathbb{R}
    \]
  \item a subset is dense if it's closure is the entire space.
    \[
      \overline{\mathcal{I}} = \mathcal{I} \cup \partial \mathcal{I}
    \]
    As demonstrated above $\partial\mathcal{I} = \mathbb{R}_u$,
    therefore, $\overline{\mathcal{I}} = \mathbb{R}_u$, and
    therefore, the closure is equal to the entire space and
    $\mathcal{I}$ is dense in $\mathbb{R}_u$
\end{itemize}
\end{document}
